\chapter{Conjugación, paridad y reversión temporal en el registro}
\label{chap:cpt-discreto-rutas-rev6}
\index{CPT!simetría discreta HMT}
\index{acción!mínima interacción}
\index{ruta!reversión temporal}

La violación CP de una matriz de mezcla y la simetría CPT no son la misma
afirmación. Antes del teorema CPT de una teoría cuántica de campos, el
registro
HMT posee una simetría discreta exacta sobre rutas. Su contenido es concreto:
conjugación intercambia las dos hojas, paridad refleja las direcciones y
reversión temporal recorre el registro en sentido contrario. La acción local
que cuenta giros y cambios de hoja es invariante bajo las tres operaciones.

\section{Ruta y acción local}

Una trayectoria de longitud \(L\) es
\[
 \gamma=
 \bigl((i_0,j_0,\ell_0),(d_1,\ldots,d_L)\bigr),
\]
donde \(d_t\in\{N,E,S,W\}\) y
\(\ell_t\in\{\Sigma,\Pi\}\). La dinámica TPK determina los estados
\((i_t,j_t,\ell_t)\). Definimos
\[
 \operatorname{turn}(\gamma)
 =
 \#\{2\leq t\leq L:d_t\neq d_{t-1}\},
\]
\[
 \operatorname{switch}(\gamma)
 =
 \#\{2\leq t\leq L:\ell_t\neq\ell_{t-1}\}.
\]

\begin{definition}[Acción discreta de mínima interacción]
\label{def:accion-minima-cpt-rev6}
Para \(\lambda\geq0\),
\[
 \mathcal A_\lambda(\gamma)
 :=
 \operatorname{turn}(\gamma)
 +\lambda\,\operatorname{switch}(\gamma).
\]
Equivalentemente,
\[
 \mathcal A_\lambda(\gamma)
 =
 \sum_{t=2}^{L}
 \left[
 1-\delta_{d_t,d_{t-1}}
 +\lambda(1-\delta_{\ell_t,\ell_{t-1}})
 \right].
\]
\end{definition}

La acción depende de la variación de la ruta, no de los nombres absolutos de
las direcciones o de las hojas.

\section{Involuciones discretas \(C\), \(P\) y \(T\) sobre las rutas: conjugación de hoja, reflexión espacial y reversión del registro}

\begin{definition}[Operadores discretos \(C,P,T\)]
\label{def:cpt-discreto-rutas-rev6}
Sobre una trayectoria completa se definen:
\begin{enumerate}
\item \(P\), reflexión espacial:
\[
 (i,j)\longmapsto(i,-j)\pmod9,
 \qquad E\leftrightarrow W,\quad N\mapsto N,\quad S\mapsto S,
\]
sin cambiar la hoja;
\item \(C\), conjugación de hoja:
\[
 \Sigma\leftrightarrow\Pi,
\]
sin cambiar la dirección; en la carta angular corresponde a
\(C^*\mapsto-C^*\), que intercambia \(q_+\leftrightarrow q_-\);
\item \(T\), reversión del registro:
\[
 d_t\longmapsto\operatorname{opp}(d_{L+1-t}),
 \qquad
 \ell_t\longmapsto\ell_{L+1-t},
 \qquad
 (i_t,j_t)\longmapsto(i_{L-t},j_{L-t}).
\]
\end{enumerate}
\end{definition}

Aquí \(T\) es una involución de la trayectoria registrada. No se identifica
con invertir un autómata empobrecido que haya olvidado el estado necesario
para reconstruir su predecesor.

\begin{theorem}[Invariancia \(CPT\) de la acción de ruta]
\label{thm:invariancia-cpt-discreta-rev6}
Para toda trayectoria \(\gamma\) y todo \(\lambda\geq0\),
\[
 \boxed{
 \mathcal A_\lambda(\gamma)
 =
 \mathcal A_\lambda(C\gamma)
 =
 \mathcal A_\lambda(P\gamma)
 =
 \mathcal A_\lambda(T\gamma)
 =
 \mathcal A_\lambda(CPT\,\gamma).}
\]
\end{theorem}

\begin{proof}
\(P\) aplica una permutación involutiva al alfabeto de direcciones. Por ello
\(d_t=d_{t-1}\) si y sólo si
\(P(d_t)=P(d_{t-1})\), y conserva el número de giros. \(C\) permuta
globalmente las hojas, de modo que conserva el patrón de igualdades
\(\ell_t=\ell_{t-1}\); no altera las direcciones. \(T\) invierte el orden de
ambas sucesiones: el número de cambios entre términos consecutivos es el
mismo al recorrer una palabra en sentido directo o inverso. Cada operación
conserva por separado \(\operatorname{turn}\) y
\(\operatorname{switch}\), luego conserva
\(\mathcal A_\lambda\). La composición también la conserva.
\end{proof}

\section{Paridad de los lectores descendentes}

El intercambio \(C^*\mapsto-C^*\) permuta \(q_+\) y \(q_-\). Por tanto,
toda expresión simétrica en los dos canales es \(C\)-par y toda diferencia
antisimétrica es \(C\)-impar. El
\cref{thm:paridad-bicapa-barbero-rev6} realiza esta regla en dos objetos:
\[
 S_{\rm bi}(-C^*)=S_{\rm bi}(C^*),
 \qquad
 \Gamma_{6\to8}(A,-C^*)
 =-\Gamma_{6\to8}(A,C^*).
\]
La banda de masa del
\cref{thm:principio-conjugacion-masa-banda} realiza la misma separación:
masa par y carga/orientación impar.

\section{Relación con CKM y con el teorema CPT de campos}

La matriz CKM construida en el capítulo siguiente posee
\(\delta_{\rm CP}\neq0\) y un invariante de Jarlskog no nulo. Eso describe
violación CP en aquella representación de mezcla. No contradice el
\cref{thm:invariancia-cpt-discreta-rev6}: este último actúa sobre el
ensemble de rutas y su acción, no exige que cada lector orientado sea par.

El teorema CPT de teoría cuántica de campos añade otro dominio y otras
hipótesis ---localidad relativista, estructura de campos, espectro y
representación correspondiente---. Se contrasta después. El resultado HMT
demostrado aquí es la simetría discreta del registro: fija el objeto que una
realización de campos debe transportar, en lugar de sustituirlo por una
referencia externa.
